% Table ASCII (codes 32 à 126), dessinée pour ce document.
\begin{center}
\footnotesize
\setlength{\tabcolsep}{3pt}
\renewcommand{\arraystretch}{1.3}
\begin{tabular}{|c|*{16}{c|}}
\hline
\textbf{hex} & \textbf{0} & \textbf{1} & \textbf{2} & \textbf{3} & \textbf{4} & \textbf{5} & \textbf{6} & \textbf{7} & \textbf{8} & \textbf{9} & \textbf{A} & \textbf{B} & \textbf{C} & \textbf{D} & \textbf{E} & \textbf{F} \\ \hline
\textbf{2} & \textsc{esp} & ! & " & \# & \$ & \% & \& & ' & ( & ) & * & + & , & - & . & / \\ \hline
\textbf{3} & 0 & 1 & 2 & 3 & 4 & 5 & 6 & 7 & 8 & 9 & : & ; & < & = & > & ? \\ \hline
\textbf{4} & @ & A & B & C & D & E & F & G & H & I & J & K & L & M & N & O \\ \hline
\textbf{5} & P & Q & R & S & T & U & V & W & X & Y & Z & [ & \textbackslash & ] & \^{} & \_ \\ \hline
\textbf{6} & ` & a & b & c & d & e & f & g & h & i & j & k & l & m & n & o \\ \hline
\textbf{7} & p & q & r & s & t & u & v & w & x & y & z & \{ & | & \} & \textasciitilde & \textsc{del} \\ \hline
\end{tabular}

\smallskip
Lecture : le caractère \texttt{A} est à la ligne \textbf{4}, colonne \textbf{1} : son code ASCII est \texttt{41} en hexadécimal ($65$ en décimal).
\end{center}
